\textcolor{red}{ TODO: check grammar and content }\\
A concept of self-replication, has steadily 
gained interest within the field of robotics as a potential paradigm for enabling 
autonomous construction, scalability, adaptability, and resilience. Self-replicating 
robotic systems (SRRS) possess the ability to autonomously build replicas of 
themselves using available components and resources. These systems hold significant 
promise for applications in remote environments such as planetary exploration,
in-situ resource utilization, and disaster recovery, where manual intervention 
is limited or infeasible \cite{Suthakorn2002SpaceUtilization} \cite{Nasa1980ConferencePublication}.

However, the complexity and autonomy involved in such 
systems pose considerable challenges to ensuring their correctness, safety, and reliability.

As SRRS evolve from theoretical constructs to practical experiments, guaranteeing 
their functional integrity becomes crucial—particularly given their recursive 
nature and potential to affect their environments in unanticipated ways. 
Traditional testing approaches are not enough when dealing with critical 
applications, vast state spaces and non-deterministic behavior inherent 
in self-replicating systems. In response to these challenges, 
formal verification methods, and in particular model checking,
have emerged as powerful tools for rigorously analyzing system behaviors against 
formal specifications.

Formal methods offer a mathematically grounded approach to verifying that a 
system adheres to desired properties such as liveness, safety, 
and correctness. 

This research explores the integration of formal verification into the design 
and analysis of self-replicating robotic systems. The goal is to bridge the 
gap between the physical nature of replication in robotics and the 
abstract frameworks used in model checking. 

The structure of this work is as follows. Section 2 provides background on formal 
verification techniques, model checking strategies, and existing self-replicating 
robotic frameworks. Section 3 outlines the core research questions for the thesis.
Section 4 describes the methodology and tools planed to use
for robotic simulations, and presents an approach for experimental validation. 

Together, these sections aim to contribute a formal 
framework for analyzing self-replicating robotic systems, supporting their future 
deployment in safety-critical and autonomous environments.
